% Figura: scenario di riferimento della GS QKD 018 (adattato da GS QKD 018, Fig. 1)
\begin{figure}[t]
\centering
\begin{tikzpicture}
  \node[box,fill=blue!8,minimum width=3.3cm] (orch) {SDN orchestrator};
  \node[box,fill=green!8,minimum width=3.3cm,below left=12mm and 6mm of orch] (cq)
       {SDN controller\\rete QKD};
  \node[box,fill=orange!10,minimum width=3.3cm,below right=12mm and 6mm of orch] (co)
       {SDN controller\\rete OTN};
  \node[box,fill=green!4,below=10mm of cq,minimum width=4.6cm] (nq)
       {nodi SD-QKD (agent, KME)};
  \node[box,fill=orange!5,below=10mm of co,minimum width=4.6cm] (no)
       {nodi OTN (SAE)};
  \draw[<->,very thick,red!70!black] (orch) -- node[lbl]{GS QKD 018} (cq);
  \draw[<->,thick] (orch) -- node[lbl]{NBI OTN (fuori scope)} (co);
  \draw[<->,thick] (cq) -- node[lbl]{GS QKD 015} (nq);
  \draw[<->,thick] (co) -- (no);
  \draw[<->,dashed] (nq) -- node[lbl]{QKD 004 / 014} (no);
\end{tikzpicture}
\caption[Scenario di riferimento della GS QKD 018]{Scenario di riferimento della GS~QKD~018 (adattato
da~\cite[Fig.~1]{qkd018}). In rosso l'interfaccia standardizzata; la
linea tratteggiata è la consegna delle chiavi dai KME ai SAE, che avviene
nello stesso perimetro di sicurezza.}
\label{fig:scenario}
\end{figure}
